\Article{4}
If we denote the relative likelihood of the total error \(x\), in a particular kind of observation, by the characteristic \(\varphi x\), then, because of the continuity of the errors, this is to be understood as meaning that the probability of an error between the infinitely close limits \(x\) and \(x+\mathrm{d}x\) is \(=\varphi x.\mathrm{d}x\). In practice it will hardly ever, indeed scarcely ever at all, be possible to assign this function a priori: nevertheless several general results concerning it can be established, which we shall present below. The following is evident. The function \(\varphi x\) must be classed among discontinuous functions insofar as, for all values of \(x\) lying outside the limits of the possible errors, it must be \(=0\); within these limits, however, it will take a positive value everywhere (leaving aside the case of which we spoke at the end of the preceding article). In most cases it will be permissible to suppose positive and negative errors of the same magnitude equally likely, in which case \(\varphi(-x)=\varphi x\). Further, since smaller errors are more readily committed than larger ones, the value of \(\varphi x\) will usually be greatest for \(x=0\), and will decrease continuously as \(x\) increases.

In general, however, the value of the integral \(\int\varphi x.\mathrm{d}x\), extended from \(x=a\) to \(x=b\), will express the probability that an error not yet known lies between the limits \(a\) and \(b\). The value of this integral from the lower limit of all possible errors to the upper limit will therefore always be \(=1\). And since \(\varphi x\)
\SourceLeaf{011}is, for all values of \(x\) outside these limits, always \(=0\), it is evident also that

\textit{the value of the integral \(\int\varphi x.\mathrm{d}x\) extended from \(x=-\infty\) to \(x=+\infty\) is always \(=1\).}

\Article{5}
Let us further consider the integral \(\int x\varphi x.\mathrm{d}x\) between the same limits, whose value we shall put \(=k\). If all the simple causes of error are such that there is no reason why one of two errors equal in magnitude but opposite in sign should be produced more readily than the other, this will also hold for the total error, or we shall have \(\varphi(-x)=\varphi x\), and consequently, necessarily, \(k=0\). From this we conclude that whenever \(k\) does not vanish, but is, for example, a positive quantity, there must necessarily be one or another cause of error that can either produce positive errors only, or at least produce positive errors more readily than negative ones. This quantity \(k\), which is in fact the mean of all possible errors, or the mean value of \(x\), can conveniently be called the constant part of the error. Moreover, it can easily be proved that the constant part of the total error is equal to the sum of the constant parts contained in the errors arising from the individual simple causes. If the quantity \(k\) is supposed known and is subtracted from every observation, and if the error of the observation thus corrected is denoted by \(x'\), and its probability by \(\varphi' x'\), we shall have \(x'=x-k\), \(\varphi' x'=\varphi x\) and consequently \(\int x'\varphi' x'.\mathrm{d}x'=\int x\varphi x.\mathrm{d}x-\int k\varphi x.\mathrm{d}x=k-k=0\), i.e. the errors of the corrected observations will have no constant part, which is also evident in itself.

\Article{6}
Just as the integral \(\int x\varphi x.\mathrm{d}x\), or the mean value of \(x\), tells us either the absence of a constant error or its presence and magnitude, the integral
\[
\int xx\varphi x.\mathrm{d}x
\]
extended from \(x=-\infty\) to \(x=+\infty\) (or the mean value of the square \(xx\)) seems best suited to defining and measuring the uncertainty of observations in general, so that, of two systems of observations differing in the likelihood of their errors, those in which the integral \(\int xx\varphi x.\mathrm{d}x\) has the smaller value are to be judged superior in precision. If anyone objects that this way has been chosen arbitrarily, without any compelling necessity, we shall readily agree.
\SourceLeaf{012}For this question, by its nature, involves something indefinite that cannot be confined within bounds except by a principle that is to some extent arbitrary. The determination of a quantity by an observation subject to a larger or smaller error may aptly be compared to a game in which there are losses only and no gains, since every error to be feared is akin to a loss. The cost of such a game is assessed by the probable loss, namely by the sum of the products of the individual possible losses and their respective probabilities. But what loss should be taken as equivalent to each error of observation is by no means evident in itself; rather, this determination depends to some extent on our choice. It is evidently not permissible to put the loss equal to the error itself; for if positive errors were taken as losses, negative ones would have to represent gains. The magnitude of the loss must instead be expressed by a function of the error that by its nature is always positive. Since the variety of such functions is infinite, the simplest that has this property seems to be preferred to the others, and this is indisputably the square: in this way the principle stated above emerges.

The illustrious \textsc{Laplace} did indeed consider the matter in a similar way, but adopted the error itself, always taken positively, as the measure of loss. But unless we are mistaken, this way is at least no less arbitrary than ours: for whether a double error is considered as tolerable as a single one repeated twice, or less tolerable, and consequently whether it is more appropriate to assign to a double error a weight only twice as great or a greater one, is a question neither evident in itself nor to be decided by mathematical demonstrations, but one to be left solely to free choice. Furthermore, it cannot be denied that continuity is violated in that way: and for this very reason that method is less amenable to analytical treatment, whereas the results to which our principle leads commend themselves by their remarkable simplicity as well as generality.

\clearpage
